\documentclass[12pt,a4paper,oneside]{article}
\usepackage[brazil]{babel}
\usepackage[utf8]{inputenc}
\usepackage{olymp}

\setcounter{problem}{4}

\begin{document}

\begin{problem}{Compactador}{compacta}{2}

Os compactadores de arquivos trabalham retirando redundância do arquivo. Por exemplo, palavras ou até partes de um texto que aparecem várias vezes podem ser escritas uma vez só, e para as demais ocorrências são gravadas apenas referências a algo que já apareceu antes. Em imagens, grandes blocos de cores iguais podem ser resumidos a uma cor e quantas vezes ela é repetida.

Um compactador bastante simples é um que troca uma sequência de caracteres iguais por um par de valores: quantas vezes o caractere ocorre, e qual é esse caractere. Veja um exemplo:

Arquivo original: \texttt{GGGGTTTTTCCAAAAAAAAAAAGGGGAAAAAACTTTTTAAAGGGGGGGGGGGGG}\\[0pt]
Compactado: \texttt{4G5T2C11A4G6A1C5T3A13G}

A descompactação também é simples, lê‐se um número e um caractere, e esse caractere é escrito repetidamente tal número de vezes.

Note que esse compactador só reduz realmente o tamanho do arquivo se houver sequência com mais de 2 caracteres iguais consecutivos. De fato, quando há apenas 2 caracteres não há ganho algum (como na primeira sequência de \texttt{C}’s, onde \texttt{CC} é substituído por \texttt{2C}). E quando não há repetição, a compactação na verdade aumenta o tamanho (como na segunda sequência de \texttt{C}, onde \texttt{C} é substituído por \texttt{1C}).

Você deve criar um programa para fazer a compactação de um arquivo como mostrado acima. O ``arquivo'', na verdade, é uma sequência de letras lidas como entrada do programa, contendo um flag no final para sinalizar o fim da sequência (um caractere ‘.’). 

%\Notes
%
%Observações, se tiver.

\InputFile

A entrada é formada por uma única linha, que contém uma sequência de letras, maiúsculas e/ou minúsculas. A linha é terminada por um caractere `\texttt{.}' (ponto), usado como flag para indicar o fim da entrada e que não deve ser compactado.

Restrição: não há limite para a quantidade total de letras mas nenhuma sequência tem mais de 1000 letras iguais seguidas.

%\Notes

\OutputFile

Seu programa deve gerar apenas uma linha na saída, contendo o resultado da compactação da entrada. Imprima ainda um `\texttt{0}' (zero) após o final da sequência compactada.

\Example

\begin{example}
\exmp{
AAAAAAAAAARRRRAAAAA.
}{
10A4R5A0
}%
\end{example}

\begin{example}
\exmp{
AAAAaaaAAAAAAa.
}{
4A3a6A1a0
}%
\end{example}

\begin{example}
\exmp{
GOOOOOOOOOOOOOOOLLLLLL.
}{
1G15O6L0
}%
\end{example}

\begin{example}
\exmp{
Pernambuco.
}{
1P1e1r1n1a1m1b1u1c1o0
}%
\end{example}

\end{problem}

\end{document}


\end{problem}

